\documentclass{article}
\usepackage[margin=2cm]{geometry}
\usepackage{amsmath}
\usepackage{graphicx}
\usepackage{booktabs}
\usepackage[utf8]{inputenc}
\usepackage{ragged2e}
\setlength{\emergencystretch}{3em}
\begin{document}
\sloppy

\noindent\hfill\includegraphics[height=1.8em]{assets/figure_header_logo.png}\par\noindent\hrule\vspace{4pt}

where NAB, SAB, and MAB are again given by, but now with the appropriate redefinition
of the coset representative as per.It should be noted that while the FDA analysis presented
in Appendix is a strong motivation for the form of the supersymmetry variations presented
above,itisnotaproof.Toactuallyderivetheform ofthesevariations,onemustfirst
introduce curvature terms representing deviations from zero of eachline in thefree differential
algebra.An application of the exterior derivative to the resulting expressions then gives rise
toBianchi identities,which must be solved before obtaining the explicitform of the fermion
variations. This is a rather involved process, and so for the moment we will content ourselves
with the motivating comments provided by the FDA. We will take the eventual presence of
smooth supersymmetric solutions consistent with the equations of motion as a posteriori
evidencefor thelegitimacyofthesevariations.Anicepropertyofthevariationsaboveis
the fact that they are consistent with the following SO(6)-invariant symplectic Majorana
condition,


\begin{equation}
\bar { \psi } _ { A } = \epsilon ^ { A B } \psi _ { B } ^ { T } \mathcal { C } \qquad ( 1 )
\end{equation}


The consistency of such a condition allows us to work with symplecticMajorana spinors just
as in the Lorentzian case, though the symplectic Majorana condition utilized here is different
than thatoftheLorentziancase.Asmentionedbefore,wewillbeconcernedwithonlythe
simplest case of a singlenon-zeroSU(2)r-charged vector multiplet scalar  3,i.e.we take
 1 =    = 0. It can be easily verified that this is a consistent truncation, and is in fact the
most general choice of non-vanishing fields that can preserve SO(4,2)   U(1)R. With this
consistent truncation, the functions NAB, SAB, and MAB appearing in the supersymmetry
variationsreduceto


\begin{equation}
\begin{array} { c } { S _ { A B } = i S _ { 0 } \epsilon _ { A B } + i S _ { 3 } \gamma ^ { 7 } \sigma _ { A B } ^ { 3 } } \\ { N _ { A B } = - N _ { 0 } \epsilon _ { A B } - N _ { 3 } \gamma ^ { 7 } \sigma _ { A B } ^ { 3 } } \\ { M _ { A B } ^ { I } = M _ { 6 } \gamma ^ { 7 } \epsilon _ { A B } + M _ { 3 } \sigma _ { A B } ^ { 3 } \qquad ( 2 ) } \end{array}
\end{equation}


wherewehavedefined


\begin{equation}
\begin{array} { r } { S _ { 0 } = \frac { 1 } { 4 } \left ( g \cos \phi ^ { 3 } e ^ { \sigma } + m e ^ { - 3 \sigma } \cosh \phi ^ { 0 } \right ) } \\ { S _ { 3 } = \frac { 1 } { 4 } i ~ m ~ e ^ { - 3 \sigma } \sinh \phi ^ { 0 } \sin \phi ^ { 3 } } \\ { N _ { 0 } = - \frac { 1 } { 4 } \left ( g \cos \phi ^ { 3 } e ^ { \sigma } - 3 m e ^ { - 3 \sigma } \cosh \phi ^ { 0 } \right ) } \\ { N _ { 3 } = - \frac {}}\end{array}
\end{equation}


Importantly, note that S3, N3,and M3 are now purely imaginary, in contrast to the Lorentzian
case.In all that follows we will set m = -1/2n such that the radius of AdSg is one.

\end{document}
